% WHAT A MATHCHOICE BRANCH MUST BEGIN WITH.
%
% \mathchoice reads FOUR sub-formulas, and each of them begins with a brace
% the reference insists on -- read the moment the branch before is done,
% rather than taken as it turns up:
%
%   $\mathchoice{a}{b}{c}\relax d$
%   ! Missing { inserted.
%   <to be read again>
%                      d
%   A left brace was mandatory here, so I've put one in.
%
% Spaces, \relax and an undefined name are stepped over on the way, the same
% as for any mandatory brace -- so the \relax above is not the branch, and
% the `d' is what the fault names.
%
% A brace put in this way opens the branch like any other, so what follows
% is swept into it and the formula's own closing token then finds a group
% still open: `Missing } inserted' for each one, once per level.
%
% HSTeX waited for a `{' to turn up and made whatever came instead into
% ordinary math material, so it drew none of this. trip line 438 writes
% \mathchoice twice, the second with only three branches before an
% \endgroup, which the reference recovers from in four faults.
\catcode`\{=1 \catcode`\}=2 \catcode`\$=3 \catcode`\^=7 \catcode`\_=8
\tracingonline=1
\font\rip=trip \fontdimen22\rip=7pt \rip
\textfont0=\rip \scriptfont0=\rip \scriptscriptfont0=\rip
\textfont1=\rip \scriptfont1=\rip \scriptscriptfont1=\rip
\textfont2=\rip \scriptfont2=\rip \scriptscriptfont2=\rip
\textfont3=\rip \scriptfont3=\rip \scriptscriptfont3=\rip
\message{[A a fourth branch that is not a brace]}
\setbox0=\hbox{$\mathchoice{a}{b}{c}\relax d$}
\message{[B a second branch that is not a brace]}
\setbox0=\hbox{$\mathchoice{a}\relax{c}{d}$}
\message{[C an endgroup where a branch belongs, with groups open]}
\setbox0=\hbox{$\begingroup{{\mathchoice{}{}{}\endgroup$}
\message{[D four branches, as they should be]}
\setbox0=\hbox{$\mathchoice{a}{b}{c}{d}$}
\message{[done]}
\end
